and the composite of these homomorphisms for @@M0@@ defines a homomorphism (7.17.4.1), which one can show (by proceeding as in 7.3) to be the same as the one defined above by means of a blow-up. The advantage of the second method, however, is that here one even has homomorphisms of the form (7.16.1):
@@M1@@
(with suitable coherent choices of the algebraic closures @@M2@@ of the @@M3@@), which allow one to define by composition a homomorphism
\sgaStableTarget{sga6.e10.eq.7-17-4-2}@@M4@@
that probably depends essentially on the choice of the @@M5@@, that is, of the @@M6@@.

\sgaStableTarget{sga6.e10.para.7-17-5}\textbf{7.17.5.} The construction of a specialization homomorphism (7.17.4.2) can be generalized to the case in which one no longer assumes that @@M7@@ is regular at @@M8@@ and that @@M9@@ is a maximal point, provided that resolution of singularities is available for the integral closed subscheme @@M10@@ of @@M11@@ with generic point @@M12@@, or merely for its localization at @@M13@@; this hypothesis holds in particular when @@M14@@ is excellent of characteristic zero, by Hironaka. Indeed, one can then find a morphism @@M15@@ with @@M16@@ regular, and, taking a closed point of the fiber @@M17@@, apply the preceding construction to @@M18@@ over @@M19@@, the generic point @@M20@@ of @@M21@@ and its closed point @@M22@@. This yields a homomorphism of the form (7.17.4.2), since @@M23@@ is a finite extension of @@M24@@ and therefore has the same algebraic closure.

\sgaStableTarget{sga6.e10.para.7-17-6}\textbf{7.17.6.} In all the constructions considered in section 7.17, the construction of a specialization homomorphism has been reduced to the
